\documentclass{article}
\usepackage{amsmath,amssymb}
\usepackage[margin=1in]{geometry}
\allowdisplaybreaks
\usepackage[UTF8]{ctex}
\begin{document}

\section*{模型 M：基于共同信息的利润分配，最终的$\alpha$表述方式}

\begin{equation}
\begin{aligned}
\max_{\alpha}\quad &
\frac{(E(P_s)/2-f-\alpha)\sigma(D)\beta}
{\left(E(D)+\dfrac{\alpha}{\sqrt{\beta-\alpha^2}}\sigma(D)\right)(\beta-\alpha^2)^{3/2}}
\left(\alpha E(D)-\sigma(D)\sqrt{\beta-\alpha^2}+\frac{E(P_sD)}{2}\right)\\
\text{subject to}\quad &
\alpha\in\left[\frac{E(P_s)}{2}-w_u,\frac{E(P_s)}{2}-w_l\right].
\end{aligned}
\tag{21}
\end{equation}
所显示的目标，实际上是该论文中对$(1-\gamma)\Pi$的表述，它代表了零售商的利润$(1-\gamma)\Pi(w(\gamma))$。关于如何最大化这一目标以及其区间范围的具体说明，则位于方程(21)之后的文字描述中，具体位置在物理PDF文件的第22页。

\paragraph{打印的终点名称及衍生标识。}
作者并未单独定义印刷后的端点名称$w_u$和$w_l$。物理PDF第22页中提到了$w(Q)\in[f,w_0]$和$\alpha=E(P_s)/2-w$。在该顺序反转的仿射映射下，该价格区间的图像为$[E(P_s)/2-w_0,E(P_s)/2-f]$。将其与(21)中的印刷区间进行比较，可以得到衍生的端点标识$w_u=w_0$和$w_l=f$。这些是源代码中支持的衍生映射，并非作者明确打印出来的别名。数量$w_0$仍然是指下面定义的全局确定的无共享响应，而不是阈值$w^{UB}$或任意的替代边界。不同的数量端点$Q_l,Q_u$，印刷后的目标和恢复公式，以及它们的端点和不确定的值，都保持如下所述。

\subsection*{固定的输入、信息和决策角色保持不变。}
$D\geq0$的需求量和$P_s\geq0$的售价都是外生的随机变量。双方都拥有相同的第一阶和第二阶矩。供应商的固定单位成本满足$f\geq0$的条件。零售商在阶段1选择利润份额，供应商在阶段2选择单位批发价格$w\geq0$，而零售商在阶段3选择产能$Q\geq0$。批发价格呈线性关系：
\[
c(Q)=wQ.
\]
未使用的容量没有回收价值。在考虑最坏情况下的分配之前，去中心化的利润以及各方的份额是
\begin{equation}
\Pi=E_{P_s,D}\!\left[P_s\min(Q,D)\right]-wQ.
\tag{2}
\end{equation}
\[
\Pi_R:=(1-\gamma)\Pi,\qquad
\Pi_S:=\gamma\Pi+(w-f)Q.
\]
这里$\gamma$是供应商的份额，$1-\gamma$是零售商的份额；(21)中的$\alpha$是转换后的批发价格变量，而不是利润份额。

给定的时刻矩阵及其分布模糊性集合是：
\[
\boldsymbol{\Sigma}=
\begin{pmatrix}
E(P_s^2)&E(P_sD)&E(P_s)\\
E(P_sD)&E(D^2)&E(D)\\
E(P_s)&E(D)&1
\end{pmatrix},
\]
\[
\Phi(\boldsymbol{\Sigma}_0)=
\left\{(P_s,D)\;;\;P_s,D\geq0,\ 
P_s,D\text{ have moments matrix }\boldsymbol{\Sigma}=\boldsymbol{\Sigma}_0\right\}.
\]
矩阵$\boldsymbol{\Sigma}_0$是固定不变的数值。这些公式中的五个矩的值都来源于该矩阵的各个元素。如第10页所述，一个非负的$P_s,D$矩矩阵是半正定的，且其各个元素都是非负的。
\[
\boldsymbol{\Sigma}\succeq0,\qquad \boldsymbol{\Sigma}\geq0.
\]
后面的不等式是逐条列出的。符号$\sigma(D)$表示需求的标准差。

\subsection*{零售商的反应以及选定的闭式形式替代方案}
原始回复以及最终表述所指向的值分别是：
\begin{equation}
Q(w):=\operatorname{argmax}_{Q\geq0}
\left\{-wQ+\inf_{(P_s,D)\in\Phi(\boldsymbol{\Sigma}_0)}
E_{P_s,D}\!\left(P_s\min(Q,D)\right)\right\},
\tag{3}
\end{equation}
\begin{equation}
\Pi(w):=\max_{Q\geq0}
\left\{-wQ+\inf_{(P_s,D)\in\Phi(\boldsymbol{\Sigma}_0)}
E_{P_s,D}\!\left(P_s\min(Q,D)\right)\right\}.
\tag{4}
\end{equation}
这些符号表示响应函数的含义；最后的$\alpha$表述采用了作者提出的闭式形式的解法，而不是额外的半定优化问题。

定义 1（第 15 页）如下：
\[
\alpha(w)=\frac{E(P_s)}{2}-w,\qquad
\beta=\frac{E(P_s^2)}{4}.
\]
定理 3 和定义 2（第 16 页）给出了阈值以及必須滿足的條件。
\[
w^{UB}=\frac12\left[
E(P_s)+
\frac{E(P_sD)E(D)-\sigma(D)\sqrt{E(P_s^2)E(D^2)-E(P_sD)^2}}
{E(D^2)}
\right],\qquad w^{UB}>f.
\]
在打印状态下：
\begin{equation}
w\leq\frac12\left[
E(P_s)+
\frac{E(P_sD)E(D)-\sigma(D)\sqrt{E(P_s^2)E(D^2)-E(P_sD)^2}}
{E(D^2)}
\right],
\tag{17}
\end{equation}
所选零售商的反应是：
\begin{equation}
Q(w)=Q^*=E(D)+\frac{\alpha(w)}{\sqrt{\beta-\alpha(w)^2}}\sigma(D).
\tag{18}
\end{equation}
根据定理 3 的替换规则，将$Q_r^*$替换为$Q^*$，那么在方程 (16) 中，对应的数值就是…
\begin{equation}
\begin{aligned}
\Pi(w)=\Pi_{wst}
&=-wQ(w)+\inf_{(P_s,D)\in\Phi(\boldsymbol{\Sigma}_0)}
E_{P_s,D}\!\left(P_s\min(D,Q(w))\right)\\
&=\alpha(w)E(D)-\sigma(D)\sqrt{\beta-\alpha(w)^2}
+\frac12E(P_sD),\qquad w\leq w^{UB}.
\end{aligned}
\tag{16}
\end{equation}
定理 3 指出，否则的话，$Q^*=0$。其证明部分（第 47-48 页）则表明……
\[
\Pi(w)=0\quad\text{for }w\geq w^{UB},\qquad
Q^*=0\quad\text{for }w>w^{UB}.
\]
特别是在$w=w^{UB}$处，印刷的闭式形式解答被保留，利润为零。第5节考虑了导致零售商订单为正的批发价格。附录B第41页指出，正响应的表达式需要$\beta>\alpha^2$；这里并未使用它来替代或缩小(21)中的印刷区间。也没有引入关于零分母或退化矩的额外约定。

\subsection*{供应商的回复被保留在最终的表述中。}
第5节限制了供应商对$f\leq w\leq w^{UB}$的回应，并进行了定义。
\[
\Pi_1(w)=(w-f)Q(w),\qquad
\Pi_S(w,\gamma)=\Pi_1(w)+\gamma\Pi(w),\quad w\leq w^{UB}.
\]
\begin{equation}
w(\gamma):=\operatorname{argmax}_{f\leq w\leq w^{UB}}
\left\{\Pi_1(w)+\gamma\Pi(w)\right\}.
\tag{19}
\end{equation}
这是所选版本所使用的供应商回复，取代了之前未经限制的$w\geq0$回复。对于$w>w^{UB}$，第5节中打印的是$\Pi_S=0$。

最终的响应分析中，将$\gamma\in[0,1]$保留其原有的顺序。同时，所有数字都保持不变。
\[
w(1)=f,\qquad w(\gamma)>f\quad\text{for all }\gamma\in[0,1).
\]
在$\gamma=1$处，零售商并不关心，因为其利润为零。尾注3中规定了强斯塔克伯格均衡原则：在多个最优决策中，零售商会选择$Q$，以最大化供应商的收益。

\paragraph{股份范围的涵盖范围。}
起始方程（5）仅打印出$\gamma\geq0$，而第5.1节至5.2节以及定理4中则使用$[0,1]$来表示份额。所选的公式为最终方程（21），其打印内容为$\alpha$区间以及端点响应规则；而之前的$\gamma$最大化目标则不再作为另一个目标来实施。

\subsection*{响应反转、端点定义以及恢复}
第 21 页上打印的反转内容如下：
\[
\alpha(Q)=
\begin{cases}
\displaystyle\frac{\sqrt{\beta}}{\sqrt{1+\left(\dfrac{\sigma(D)}{Q-E(D)}\right)^2}},&Q\geq E(D),\\[6pt]
\displaystyle-\frac{\sqrt{\beta}}{\sqrt{1+\left(\dfrac{\sigma(D)}{Q-E(D)}\right)^2}},&Q<E(D),
\end{cases}
\]
\[
w(Q)=\frac{E(P_s)}2-\alpha(Q)=
\begin{cases}
\displaystyle\frac{E(P_s)}2-\frac{\sqrt{\beta}}{\sqrt{1+\left(\dfrac{\sigma(D)}{Q-E(D)}\right)^2}},&Q\geq E(D),\\[6pt]
\displaystyle\frac{E(P_s)}2+\frac{\sqrt{\beta}}{\sqrt{1+\left(\dfrac{\sigma(D)}{Q-E(D)}\right)^2}},&Q<E(D).
\end{cases}
\]
定义 3 及其相关文本如下：
\[
w_0=w(\gamma=0),\qquad f<w_0\leq w^{UB},\qquad
Q_u=Q(f),\qquad Q_l=Q(w_0),
\]
\[
w(Q)\in[f,w_0],\qquad Q\in[Q_l,Q_u].
\]
对于内部响应，该论文采用了一阶条件。
\[
\frac{\partial\Pi_S}{\partial w}(w^*,\gamma)
=\frac{d\Pi_1}{dw}(w^*)+\gamma\frac{d\Pi}{dw}(w^*)=0
\]
并且声明
\begin{equation}
\gamma(Q)=-\frac{\dfrac{d\Pi_1}{dw}(w(Q))}
{\dfrac{d\Pi}{dw}(w(Q))}.
\tag{20}
\end{equation}
附录 B 中，第 48 至 49 页所使用的导数定义如下：
\begin{align}
\frac{dQ}{dw}(w)&=-\frac{\beta\sigma(D)}{(\beta-\alpha^2)^{3/2}}<0,
&&f\leq w\leq w^{UB},\tag{48}\\
\frac{d\Pi}{dw}(w)&=-\left(E(D)+\sigma(D)\frac{\alpha}{\sqrt{\beta-\alpha^2}}\right)=-Q(w),\tag{49}\\
\frac{d\Pi_1}{dw}(w)&=Q(w)+(w-f)\frac{dQ}{dw}(w),\tag{50}\\
\frac{\partial\Pi_S}{\partial w}(w,\gamma)&=(1-\gamma)Q(w)+(w-f)\frac{dQ}{dw}(w),\tag{51}
\end{align}
在这些派生公式中，$\alpha=\alpha(w)$这样的数学符号应保持不变，按照原来的顺序进行复制。所有数字都需保留不变。

定理 4 指出，当零售商做出选择时，结果就是 $Q\in[Q_l,Q_u]$。
\[
\gamma=1-
\frac{(E(P_s)/2-f-\alpha(Q))\sigma(D)\beta}
{(\beta-\alpha(Q)^2)^{3/2}Q},\qquad Q\in(Q_l,Q_u),
\]
带有单独端点选择的
\[
Q=Q_l:\ \gamma=0,\qquad Q=Q_u:\ \gamma=1.
\]
文本中还指出，$\gamma(Q)=0$成立的条件是$w^*=w_0$成立，而$\gamma(Q)=1$成立的条件是$w^*=f$成立。

对于优化器$\alpha^*$（21），打印出的恢复方程如下：
\begin{equation}
\gamma^*=1-
\frac{(E(P_s)/2-f-\alpha^*)\sigma(D)\beta}
{\left(E(D)+\dfrac{\alpha^*}{\sqrt{\beta-(\alpha^*)^2}}\sigma(D)\right)
(\beta-(\alpha^*)^2)^{3/2}}.
\tag{22}
\end{equation}

\paragraph{保留了来自原材料的边界和充分性条件。}
印刷版中的分支$Q\geq E(D)$包含$Q=E(D)$，尽管其表达式被$Q-E(D)$除去了；在该等式中没有提供任何值，因此这里也没有添加任何内容。第21页明确指出一阶条件是必要的，但可能是不够的，并声称(20)中的份额能够诱导出期望的供应商响应。定理4保留了这一论断，并分别给出了端点选择。目标函数(21)和恢复函数(22)在封闭的$\alpha$区间内被列出，没有单独的端点限定，而定理4中的明确内部份额公式则限定在$Q\in(Q_l,Q_u)$范围内。这些范围被保留下来，而不是被修正。没有使用任何不可用的外部补充材料或勘误表；来自这些来源的任何更正都未经验证。

\end{document}
